\documentclass[conference,10pt]{IEEEtran}
\IEEEoverridecommandlockouts
\usepackage{graphicx,booktabs,amsmath,amssymb,xcolor,url,xspace}
\usepackage[hidelinks]{hyperref}
\hypersetup{pdftitle={로봇 양팔을 위한 에고센트릭 인간 시연: 방법과 파이프라인}, pdfauthor={Jeong Hwan Lee}}
\usepackage{tikz}
\usetikzlibrary{positioning,arrows.meta}
\usepackage{fontspec}\setmainfont{Times New Roman}\usepackage{xeCJK}\setCJKmainfont[AutoFakeBold=2.5]{AppleMyungjo}\setCJKsansfont[AutoFakeBold=2.5]{Apple SD Gothic Neo}\setCJKmonofont{Apple SD Gothic Neo}\xeCJKsetup{CJKecglue={},CJKspace=true}
\newcommand{\mm}{\,mm\xspace}
\def\UrlBreaks{\do\/\do\_\do\.\do\-}\def\UrlFont{\small\ttfamily}\newcommand{\code}[1]{\path{#1}}

\begin{document}
\title{로봇 양팔을 위한 에고센트릭 인간 시연\\방법과 파이프라인}
\author{\IEEEauthorblockN{Jeong Hwan Lee}
\IEEEauthorblockA{기술 보고서, 버전 4, 2026년 10월 7일 \quad alnrg2arg@gmail.com\\
제1부(큐브 삼단 쌓기)와 제2부(단일 팔 접근)의 동반 문서. 각 단계의 결과가 아니라 작동 방식을 기술한다.}}
\maketitle

\begin{abstract}
본 문서는 에고센트릭 인간 시연을 로봇 정책으로 바꾸는 데 사용한 모든 단계를 알고리즘, 파라미터, 게이트 수준에서 기술한다. 손에
착용하는 기록 장치와 수집 프로토콜, 센서와 로봇의 캘리브레이션, 동기화와 내보내기, 단안 손목 자세 복원과 세 가지 메트릭 스케일
소스(IMU, 크기를 아는 큐브, 테이블 평면), 로봇 학습 데이터로의 두 세대 변환(관절 공간 리타게팅과 데카르트 액션 계약), 작업
프레임을 공유하는 단일 팔 변형, X-VLA 모델과 학습 레시피, 오프라인 평가와 진단 평가, 그리고 로봇 배포를 다룬다. 별도 언급이 없으면
파일 이름은 저장소 기준이다. 마지막 절에는 이 문서를 쓰는 동안 우리 코드에서 발견한 불일치를 나열한다.
\end{abstract}

\section{개요}
그림~\ref{fig:overview}은 각 단계를 보여 준다. 각 단계는 출력을 매니페스트와 함께 기록하며, 그 뒤에 오는 게이트가 단위(에피소드,
세그먼트, 행 또는 테이크)를 승인하거나 사유와 함께 거부한다. 보고서는 각 게이트가 제거한 양을 센다. 두 변환은 앞단을 공유한다.
1세대(Gen-1)는 인간 손목 자세를 로봇 관절 궤적으로 리타게팅하고, 2세대(Gen-2)는 이를 우리 로봇 원격조작 데이터와 같은 상대
데카르트 액션 공간으로 표현한다. 단일 팔 접근 데이터(제2부)는 일부 차원을 마스킹하고 로봇과 물리적 원점을 공유하는 Gen-2 변형이다.

\begin{figure*}[t]
\centering
\begin{tikzpicture}[font=\scriptsize, node distance=3mm and 3mm,
  box/.style={draw=black!55, rounded corners=1.5pt, align=center, inner sep=2pt, minimum height=9mm, text width=24mm, fill=white},
  g1/.style={box, fill=orange!8, draw=orange!60!black}, g2/.style={box, fill=blue!7, draw=blue!55!black},
  dep/.style={box, fill=green!7, draw=green!45!black}, arr/.style={-{Stealth[length=1.5mm]}, draw=black!60}]
\node[box] (rig) {\textbf{장치 + 프로토콜}\\2$\times$HandUMI, 헤드 C922, AUTO 루프};
\node[box, right=of rig] (cal) {\textbf{캘리브레이션}\\어안, 카메라--IMU, 카메라$\to$TCP, 조, 로봇 카메라};
\node[box, right=of cal] (exp) {\textbf{내보내기 + 동기화}\\시계 피팅, $\le$20\,ms\\게이트, 그리퍼 폭};
\node[box, right=of exp] (pose) {\textbf{손목 자세}\\MASt3R-SLAM (RGB), 점프 필터};
\node[box, right=of pose] (scale) {\textbf{메트릭 스케일}\\IMU-VI $\mid$ 큐브 PnP $\mid$ 원점 평면};
\node[g1, below=6mm of rig] (g1) {\textbf{Gen-1}\\IK 리타게팅 $\to$ 의사 관절,\\65프레임 세그먼트};
\node[g2, right=of g1] (g2) {\textbf{Gen-2 CART20}\\상대 자세 16개, 15\,Hz 행, IK 필터 없음};
\node[g2, right=of g2] (hra) {\textbf{단일 팔}\\loss mask, G 프레임,\\테이블 인지 QP};
\node[box, right=of hra] (train) {\textbf{X-VLA 0.9B}\\사전학습 / 파인튜닝 /\\공동 학습 (Ray)};
\node[dep, right=of train] (dep) {\textbf{로봇}\\MPS 추론, Pink IK, PLAN 실행};
\draw[arr] (rig)--(cal); \draw[arr] (cal)--(exp); \draw[arr] (exp)--(pose); \draw[arr] (pose)--(scale);
\draw[arr] (scale.south) |- ++(0,-2.6mm) -| (g1.north); \draw[arr] (scale.south) |- ++(0,-2.6mm) -| (g2.north);
\draw[arr] (g2)--(hra); \draw[arr] (hra)--(train); \draw[arr] (train)--(dep);
\draw[arr] (g1.south) |- ++(0,-2.5mm) -| (train.south);
\end{tikzpicture}
\caption{파이프라인 개요. 주황: 1세대(관절 공간). 파랑: 2세대(데카르트)와 그 단일 팔 변형. 초록: 배포. 모든 화살표 뒤에는 사유를
기록하는 게이트가 있다.}
\label{fig:overview}
\end{figure*}

\section{기록 장치와 수집 프로토콜}\label{sec:rig}
\textbf{장치} (\code{collection/configs/handumi/hardware_handumi_v1.yaml}). 각 HandUMI 유닛에는 Arducam 어안 UVC 카메라
(1920$\times$1080, 30\,fps, MJPG), Teensy~4.1에 연결된 ICM-42688-P IMU(공칭 200\,Hz, 보드마다 자체 오실레이터를 가지므로 실측은
201.6, 197.9\,Hz), 조(jaw) 인코더로 쓰는 Feetech STS3215 서보(4096 ticks/rev, 100\,Hz 폴링)가 달려 있다. 두 손목 카메라가 같은
USB 시리얼을 보고하므로 제품 이름으로 구분한다. 고정된 Logitech C922(640$\times$480, 30\,fps, 자동 초점 끔)가 탁자를 기록한다.
RGB-D 카메라는 손목 스트림을 방해하여 제외하였다.

\textbf{기록} (\code{collector/recorder.py}). 에피소드 디렉터리에는 카메라별 MP4 하나, \code{sensors.mcap}, \code{events.json},
\code{episode_meta.json}이 들어 있으며, 닫힐 때에만 이름이 \code{.incomplete}에서 \code{.complete}로 바뀐다. MCAP JSON 채널에는
IMU 샘플(장치·호스트 타임스탬프, 가속도계, 자이로), 그리퍼 샘플(언랩·랩 tick, 정규화 값), 프레임별 메타데이터(프레임 인덱스,
캡처 시각, 건너뛴 프레임), 이벤트, 시계 앵커가 담긴다. 모든 스트림에는 호스트 단조 시계로 타임스탬프를 찍는다. IMU는 호스트 수신
시각, 그리퍼는 시리얼 트랜잭션의 중간 시점, 프레임은 캡처 시각을 쓴다.

\textbf{AUTO 루프} (\code{collector/autoloop.py}). RESET $\to$ ARMING $\to$ HOLD $\to$ (PREP) $\to$ EPISODE $\to$ STOP 상태
기계가 음성 신호와 함께 각 테이크를 진행한다. ARMING은 정지 상태를 기다린다. 0.25\,s 창에서 자이로 크기의 평균이 6\textdegree/s
미만, 가속도 크기의 표준편차가 0.35\,m/s$^2$ 미만, 최대 샘플 간격이 25\,ms 미만, 창 채움률이 60\% 이상인 상태가 1\,s 동안 연속되어야
한다. 기록 시작 후의 정지 구간은 시각--관성 초기화 창 역할을 한다. 예비 QA를 통과한 테이크는 자동으로 보존되고, 실패하면 루프가
일시 정지한다.
\begin{itemize}
\item \emph{쌓기}: 20\,s 테이크, 10\,s 리셋, 3\,s 정지. 다음 순서는 보존된 테이크가 가장 적은 순서로 정해지므로
(\code{episode_manager.next_order}) 순서는 무작위가 아니라 균형을 맞춘 것이다.
\item \emph{로봇 유사}(\code{robot_like_v1}): 실시간 피드백이 있는 30\,s 테이크. 15\,Hz 격자에서 손목 각속도와 각가속도를 로봇
원격조작 데이터의 95번째 백분위수(0.996\,rad/s, 4.413\,rad/s$^2$)와 비교하며, 이를 넘는 틱이 10\% 초과이거나 파지가 두 번 미만이면
테이크를 검토로 보낸다.
\item \emph{접근}(HRA\_red): 오른손만, 10\,s 테이크, 5\,s 리셋, 그리퍼 검사 끔.
\item \emph{원점이 있는 접근}(HRA\_A100): 5\,s 휴식, 조 끝을 테이프로 표시한 십자에 대고 2\,s 정지, 2\,s 원점 정지를 추가로 하며
기록 시작(이벤트 \code{go_cue}), 로봇 유사 시작 자세로 3\,s ``ready'' 이동, ``go''에서 10\,s 접근. 오른 손목은 로봇 카메라와 비슷한
뷰도 기록한다.
\end{itemize}

\textbf{수집 QA} (\code{collector/integrity.py}). 장치나 레코더 오류 이벤트, 빈 스트림, 25\,fps 미만의 프레임률이 있으면 테이크를
검토로 보낸다. 그리퍼가 20\,Hz 미만으로 샘플링되거나, 멈춰 있거나, 캘리브레이션 범위를 300 tick 넘게 벗어나거나, 이동 거리가 범위의
5\% 미만이면 거부한다. 검증 단계에서는 모든 비디오를 디코딩하여 메타데이터 프레임 수와 맞는지 확인하고, 프레임 인덱스가 연속이며
캡처 시각이 엄격히 증가하는지 검사한다.

\section{캘리브레이션}\label{sec:cal}
\textbf{손목 카메라.} 전체 해상도에서 Kannala--Brandt(등거리) 어안 모델을 사용하며, 25\mm{} 정사각형의 8$\times$5 체커보드로
Kalibr로 피팅한다(왼쪽 $f_x{=}788.9$\,px). \textbf{카메라--IMU.} 손별 Kalibr~\cite{furgale2013kalibr} 외부 파라미터와 시간
오프셋. \textbf{카메라에서 TCP로} (\code{calibration/camera_tcp/handumi_camera_tcp_v2.yaml}). TCP는 집게 중심이며, $x$는 접근
방향, $y$는 조가 닫히는 방향, $z = x\times y$이다. 회전은 고정된 축 순열($x_\text{TCP}=z_\text{cam}$, $y_\text{TCP}=-x_\text{cam}$,
$z_\text{TCP}=-y_\text{cam}$)이고, 병진은 캘리퍼 측정값(앞으로 95\mm, 아래로 95\mm)을 IMU 중력으로 측정한 카메라 피치
(20.7\textdegree)만큼 회전한 것으로, $t=(0, 55.3, 122.5)$\mm{}를 얻으며 큐브 기반 피팅과 약 1\mm{} 이내로 일치한다. 이 값은
유닛 하나에서 측정하여 양쪽에 적용하였다. \textbf{헤드 카메라.} 임시 plumb-bob 모델($f=677.7$\,px).

\textbf{로봇 측}(제2부). 로봇의 고정 카메라는 같은 체커보드로 캘리브레이션한다($f=728.7$\,px, RMS 1.05\,px). 오른쪽 조 끝으로
내부 코너 네 개를 터치하고, 순기구학(FK) TCP 위치와 PnP로 구한 보드 자세를 Kabsch 정렬하여 $T_\text{base,board}$를 얻으며,
$T_\text{base,cam}=T_\text{base,board}T_\text{cam,board}^{-1}$이다. 로봇 손목 카메라는 보드 뷰 13장으로 캘리브레이션하고(종횡비
고정, 접선 왜곡 없음, RMS 0.16\,px), 핸드-아이 변환은 Tsai, Park, Horaud, Daniilidis 방법~\cite{tsai1989handeye}으로 풀어 Park의
해를 사용한다. 물리적 접촉은 URDF 손가락 끝 박스(IK TCP에서 $x$를 따라 뒤로 70--80\mm, $z$로 $\pm$19.6\mm)로 모델링하며, 그
가장 낮은 모서리를 접촉점으로 간주한다.

\section{내보내기와 동기화}\label{sec:export}
\textbf{시계 정렬} (\code{pose/timing.py}). USB 지연은 항상 지연을 더하기만 하므로, IMU 장치 시각을 최소제곱 기울기로 호스트 시각에
대응시키고 절편은 잔차의 5번째 백분위수에 다시 고정한다. 그다음 캘리브레이션에서 얻은 카메라--IMU 시간 오프셋을 적용한다. 카메라는
조금씩 다른 속도(30.003--30.026\,Hz)로 자유 동작하므로 프레임은 고정 오프셋이 아니라 가장 가까운 캡처 시각으로 짝짓는다.

\textbf{내보내기} (\code{pipeline/export/umi_export_orbslam.py}). 각 손목은 960$\times$540 비디오(내부 파라미터 절반), 첫 프레임
이전 50\,ms의 IMU를 포함한 GoPro 텔레메트리 형식의 IMU 파일, 카메라 설정 파일로 기록된다. 프레임과 캡처 시각은 단일 디코딩 패스에서
얻으므로 서로 어긋날 수 없다.

\textbf{동기화 게이트} (\code{pipeline/census/c8_census_pre.py}, \code{c8_phase3.py}). 왼 손목이 마스터 시계이고, 오른 손목과 헤드
카메라는 가장 가까운 캡처 시각으로 매칭한다. 두 손목 자세가 모두 유효하고 두 매칭이 모두 20\,ms 이내($\Delta t_R<20$\,ms,
$\Delta t_H<20$\,ms)이면 행이 유효하다. 연속된 유효 행 65개 이상의 구간이 하나라도 있으면 에피소드가 통과한다.

\textbf{그리퍼 폭} (\code{pipeline/grip/}). 각 프레임은 가장 가까운 그리퍼 샘플을 취한다. 기록된 정규화 값은 신뢰할 수 없었으므로
(열린 끝이 기계적 정지점 너머에 있었고, 왼쪽 방향이 세션마다 뒤집혔다) Gen-1은 raw tick으로부터 폭을 다시 계산한다. 방향은 열림/닫힘
tick 쌍에서, 닫힘 앵커는 에피소드별 닫힘 극값에서 정하고, 열린 끝은 닫힘 앵커에 장치 범위(1136.6 / 1248.2 ticks)를 더한 값이며,
폭 $= \mathrm{clip}(f,0,1)\times 67.5$\mm{}이다. Gen-2는 이후의 캘리퍼 캘리브레이션(gcal1)을 사용한다. 0/20/40/60/80\mm{}에 구간별
선형 보간점을 두고 $g=\mathrm{clip}(w/80\,\text{mm},0,1)$로 정규화하며, 0 = 닫힘이다.

\section{손목 자세와 메트릭 스케일}\label{sec:pose}
\textbf{단안 자세} (\code{pipeline/mast3r_pose/run_perframe.py}). MASt3R-SLAM~\cite{murai2025mast3rslam}(공식 코드, 커밋 고정,
유일한 패치는 빌드 수정)을 핀홀 뷰로 왜곡 보정한 RGB 손목 비디오에 대해 단일 스레드, 매 프레임, 고정 시드로 실행한다. 전역 최적화
후 모든 프레임을 최종 키프레임 기준으로 다시 표현한다.
\begin{equation}
T_{WC}(f) = T_{WK}^{\text{final}}(k)\,T_{WK}^{\text{track}}(k)^{-1}\,T_{WC}^{\text{track}}(f),
\end{equation}
이는 Sim(3) 사슬이다. 재위치추정(relocalisation) 중에 추적된 프레임은 이후 키프레임이 되지 않는 한 무효이다. 장면당 ORB-SLAM3 맵
하나를 쓰는 UMI 방식~\cite{chi2024umi}은 장면이 바뀐 뒤 재위치추정에 실패하여 대체하였다.

\textbf{IMU 시각--관성 스케일} (\code{cp6_scale.py:imu_vi}). SLAM 궤적은 에피소드별 스케일 $s$를 제외하면 메트릭이다.
$R_{wb}=R_{wc}R_{cb}$, 각 프레임 주위에서 박스 평균한 가속도계 샘플 $f_b$, 그리고 SLAM 카메라 위치와 $R_{wc}t_{cb}$의
Savitzky--Golay 2차 도함수(차수 3, 창 $W$)로 얻은 카메라 가속도 $a_\text{cam}$과 레버암 가속도 $a_\text{lever}$를 사용하면, 각
프레임은 다음 식의 세 행을 제공한다.
\begin{equation}
R_{wb}\,f_b - a_\text{lever} = s\,a_\text{cam} - g\;\;(+\,R_{wb}b_a),
\end{equation}
이를 가장 긴 유효 구간에 대해 Huber 가중 반복 재가중 최소제곱(10회 반복, $\sigma=1.4826\,$MAD)으로 $s$와 $g\in\mathbb{R}^3$에 대해
푼다. 양변이 같은 필터를 거치므로 대역 제한이 같고, 어떤 것도 두 번 적분하지 않는다. 바이어스 없는 $W=15$는 마커 벤치마크의 피팅
분할에서 선택하였다. $s>0$, $|g|\in[8.81,10.81]$\,m/s$^2$이고 $s$가 벤치마크 최소 스케일의 0.1$\times$와 최대 스케일의 10$\times$
사이에 있으면 스케일이 유효하다. 메트릭 TCP 자세는 $T_\text{tcp}=(T_{WC}\text{ with translation}\times s)\,X_\text{cam,tcp}$이다.

\textbf{자세 벤치마크.} 80\mm{} ArUco 마커 네 개가 보이는 시연 14개(손 트랙 28개, 피팅 9개, held-out 5개). 마커 자세는 왜곡 보정한
점에 대해 IPPE로 풀며, 합의 방향에서 20\textdegree{} 넘게 벗어난 마커는 버린다. 정답 스케일은 열 프레임 떨어진 프레임 쌍 중 마커가
30\mm{} 넘게 움직인 쌍에서 마커 변위와 SLAM 변위의 비의 중앙값이다. 커버리지, 스케일 오차, 회전 오차, UMI 방식의 16스텝 상대 위치
오차, 방향 반전 비율을 보고한다.

\textbf{카메라--IMU 일관성.} 자이로와 SLAM 각속도가 일치하지 않는 세션(다른 곳의 0.92--0.98 대비 코사인 중앙값이 0 근처)은 격리한다.
접근 데이터에서는 큐브 및 테이블 기하와 비교한 뒤, 오른 손목의 내보낸 카메라--IMU 회전을 카메라 $x$축에 대해 약
18\textdegree{} 수정하였다.

\textbf{자세 점프 필터} (\code{cart20/ego_cart20/preprocessing/pose_filter.py}). MASt3R-SLAM은 가끔 손실 플래그 없이 한 프레임 안에
카메라를 수 센티미터 움직인다. 한 66.7\,ms 행 동안의 변위가 100\mm{} 또는 45\textdegree{}를 넘고 점프가 그 스텝 안에 있을 때, 또는
50\mm{}를 넘는 스텝이 0.3\,m/s보다 느린 스텝들 사이에 고립되어 있을 때 raw 스텝에 플래그를 단다. 5샘플 이내의 플래그는 하나의
이벤트로 묶는다. 자세가 되돌아오면($\max(30\,\text{mm}, 0.3\times\text{step})$와 45\textdegree{} 이내) 이벤트는 스파이크이며 해당
샘플만 무효화한다. 그렇지 않으면 단방향 점프로 보고 그 팔 트랙의 나머지를 무효화한다. 시작점에 고정된 상태가 그 이후로 틀리게 되기
때문이다.

\textbf{큐브 PnP 스케일} (\code{cart20/ego_cart20/cube_pnp.py}, 제2부). 느린 동작에서는 IMU 피팅이 관측 불가능하다. 빨간 큐브(모서리
38\mm)를 HSV로 분할하고, 볼록 껍질을 꼭짓점 여섯 개(또는 네 개)로 줄인 뒤 각 변의 가운데 부분에 대한 직선 피팅으로 모서리를
날카롭게 다듬는다. 이를 2면·3면 실루엣 모델로 PnP 피팅하거나(가시성 검사: 카메라가 앞쪽에 놓이는 면들이 모델의 보이는 면과 같아야
한다), 정면에서 35\textdegree{} 이내이면 중심을 면 뒤쪽 반 모서리 거리에 둔 단일 면 IPPE 모델로 피팅한다. 큐브가 정지해 있으므로 그
중심 $c_i=p_i + k\,R_i t_i$는 일정해야 하며, 이로부터 역스케일을 닫힌 형태로 얻는다.
\begin{equation}
k = -\frac{\sum_i (p_i-\bar p)\cdot(u_i-\bar u)}{\sum_i\lVert u_i-\bar u\rVert^2},\qquad u_i=R_i t_i,\quad s=1/k,
\end{equation}
여기에 MAD 트리밍과 부트스트랩 재표본 200회를 적용한다. 테이크는 PnP 프레임 10개 이상, 부트스트랩 퍼짐 10\% 이하, 중심 잔차
2\,cm 이하를 만족해야 한다.

\textbf{원점 평면 스케일}(제2부, HRA\_A100). 원점 정지 동안 손가락 끝이 판에 닿아 있으므로 첫 SLAM 키프레임은 알려진 거리에서
테이블을 본다. 손가락 영역과 신뢰도가 낮은 점을 제외하고 RANSAC(400회 반복, 임계값은 깊이 중앙값의 0.4\%)으로 평면을 최대 세 개
피팅하며, 첫 평면과 평행한 것 중 가장 큰 평면을 테이블로 택한다. 그러면 $s=|n\cdot t_\text{tip}|/(s_0\,|d|)$이며, $t_\text{tip}$은
카메라 좌표계의 손가락 끝, $s_0$는 키프레임의 Sim(3) 스케일이다. 큐브 PnP가 유효하면 그것을, 아니면 평면을 사용한다.

\section{1세대: 관절 공간 리타게팅}\label{sec:gen1}
\textbf{세그먼트와 게이트} (\code{c8_phase3.py}). 유효 행을 겹치지 않는 65프레임 블록(2.17\,s)으로 자른다. 팔별로 세그먼트는 유효한
메트릭 스케일을 가져야 하고(B), 행의 95\% 이상이 로봇 작업 공간 안에 있어야 하며(C; 목표는 로봇 원격조작 TCP 자세 6{,}000개 중
최근접 이웃이 50\mm{} 이내이고 가장 가까운 50개의 방향 차이 중 최솟값이 20\textdegree{} 미만이면 내부로 본다), 모든 행에서 IK가
풀려야 하고(D), 유한한 관절 값을 내야 한다(E). 두 팔이 모두 통과하면 세그먼트는 양팔 세그먼트이다.

\textbf{리타게팅} (\code{pipeline/pseudo_joint/pseudo_joint_pipeline.py}). 인간 동작을 세그먼트 시작 기준 상대 동작으로 취해 로봇
데이터셋의 TCP 프레임으로 옮기고($\text{rel}=F^{-1}T_{0}^{-1}T_{t}F$, $F=R_x(180^\circ)$), 고정된 로봇 앵커 자세(로봇 원격조작
데이터의 활동 자세 중앙값)에 배치한다. $T_\text{tgt}=\mathrm{FK}(q_\text{anchor})\,C^{-1}\,\text{rel}\,C$, $C=R_y(90^\circ)$.
최소제곱 IK(방향 가중치 20의 자세 비용, 이전 해를 향한 rest 항, 관절 한계 제약, 최대 10회 반복, 관절 스텝 최대 0.35\,rad)는 이전
행에서 시드를 받는다. 행에는 LIMIT, DQ\_FAIL($>$1\,rad), BRANCH\_SWITCH(목표가 거의 움직이지 않는데 $>$0.3\,rad), NOT\_CONVERGED
($>$15\mm{} 또는 $>$5\textdegree), OK 중 하나의 레이블을 붙인다.

\textbf{데이터셋과 목표} (\code{pipeline/dataset/}, \code{training/humanik_delta.py}). 세그먼트당 LeRobot 에피소드 하나이며,
640$\times$480 뷰 세 개, 14차원 상태(팔마다 관절 각도 여섯 개와 그리퍼 값 하나), 목표 TCP 자세로 구성된다. 정책은 누적 관절 변위
30개 $a_i=q_{t+5+i}-q_t$(5프레임 선행)와 절대 이진 그리퍼를 예측한다. 두 보조 손실이 관절 공간을 측정된 동작에 묶는다. (C) 여분
액션 차원에 상대 TCP 병진 $\mathrm{trans}(T_\text{tgt}(t)^{-1}T_\text{tgt}(t{+}5{+}i))\times10$을 두는 항과, (D) FK 일관성 항이다.
\begin{equation}
\mathcal{L}=\mathcal{L}_{\Delta q}+\mathcal{L}_\text{grip}+2.0\,\mathcal{L}_C+20\,\lVert \mathrm{FK}(q_t+\hat a)-p_\text{tgt}\rVert^2 .
\end{equation}
분할은 고정 시드로 원본 에피소드의 10\%를 뽑으며, 양쪽에 걸치는 에피소드는 없다.

\textbf{리타게팅 v2} (\code{pipeline/retarget_v2/}). 단일 앵커 때문에 모든 세그먼트가 한 자세에서 시작했다. v2는 캘리브레이션에 로봇
에피소드 30개만 사용한다(다른 120개로 평가). 두 팔 관절의 $2^{17}$개 Sobol 샘플로 된 시드 뱅크를 작업 공간, 테이블 바닥, 조작성
(위치 야코비안의 최소 특이값), 팔 간 간격 $\ge$75\mm{}(팔당 링크 세그먼트 9개 사이의 정확한 세그먼트 간 거리)로 거른 뒤 중복을
제거한다(시드 1{,}195개). 변형 K는 상위 시드 8개를 롤아웃하고
$S=\overline{e_\text{pos}}/15+\overline{e_\text{rot}}/5+\overline{|\Delta q|}/0.6+0.5\,\overline{|\Delta^2q|}/0.6$로 점수를 매긴다.
Kr은 로봇 데이터까지의 자세 거리를 추가한다. TR은 96차원 상대 동작 특징으로 가장 가까운 로봇 궤적 창 64개를 검색하고, 실현 가능한
것을 남겨 $q_\text{last}+\Delta q_\text{ref}$로 IK 시드를 준다. 하이브리드는 TR이 실현 가능하면 TR을, 아니면 Kr을 쓴다.

\section{2세대: 데카르트 계약(CART20)}\label{sec:gen2}
\textbf{계약} (\code{cart20/ego_cart20/config.py}). 행은 15\,Hz이고, 오프셋 $\Delta=3/59.94$\,s $=50.05$\,ms는 로봇 데이터의 스텝이다.
행 시각 $t$와 각 팔에 대해
\begin{equation}
A_k = T(t)^{-1}\,T(t+k\Delta),\qquad k=1,\dots,16 ,
\end{equation}
이며, 순차가 아니라 현재 시점에 고정되고 raw 30\,Hz 트랙에서 보간한다(위치는 선형, 회전은 SLERP, 그리퍼는 선형; 50\,ms 이내로
떨어진 유효 샘플 사이에서만; 외삽 없음). 각 팔은 [위치 3 $\mid$ 6차원 회전($R$의 처음 두 행, Gram--Schmidt로 디코딩) $\mid$
그리퍼]를 내므로 두 팔은 20차원이 되며, 모델의 32차원 액션은 정확히 0인 12차원으로 패딩된다. 상태는 각 팔의 자체 작업 시작 자세
기준 상대 자세 $T(t_0)^{-1}T(t)$와 두 그리퍼 값이다(좌우가 공유하는 SLAM 프레임은 없다).

\textbf{변환} (\code{convert.py}). 자세 점프 필터, 팔별 보간 트랙, 두 팔이 모두 유효한 첫 순간을 작업 시작으로, 20\,ms 이내의 가장
가까운 헤드 프레임. 시각 $t$에서 두 팔이 모두 유효하고, 목표 16개가 모두 유효하며, 헤드 프레임이 있을 때에만 학습 행이 되고 패딩은
없다. 작업 공간이나 IK에 의한 선별은 적용하지 않는다.

\textbf{검증} (\code{validation/}, \code{tests/}). 형상, 0 패딩, 유한 값, $g\in[0,1]$, 행렬식 $+1$인 정규직교 회전, 항등 시작 상태,
서로소 분할, 마지막 유효 행을 넘는 목표 없음, $k$에 따라 증가하는 병진 중앙값을 검사하며, 합성 테스트 11개가 계약을 고정한다.

\textbf{로봇 유사 손목 렌더링} (\code{scripts/export_lerobot_robotized.py}). 시각적 격차를 줄이기 위해 어안 프레임을 같은 방향의
가상 핀홀 카메라로 재렌더링한다(따라서 레이블은 그대로 유효하다). 수평 화각은 62--70\textdegree{}에서 뽑고, 로봇처럼 조가 아래쪽에
오도록 주점을 75--125\,px 아래로 옮긴 뒤 축소, 블러($\sigma$ 0.4--1.0), JPEG 압축(품질 55--85), 광도 지터를 프레임별 시드로
적용한다. ROBOT100은 이를 모든 프레임에 적용한다. 단일 팔 robotcam v2 렌더링(\ref{sec:hra}절)은 대신 측정된 로봇 카메라를 사용한다.

\textbf{내보내기.} 15\,fps의 LeRobot v3, 224$\times$224 뷰 세 개, 통계 재계산(패딩 차원은 평균 0, 표준편차 1). 관절 필드는 일부러
NaN으로 두어 관절 공간 손실이 쓰이면 바로 눈에 띄게 실패하도록 한다.

\section{단일 팔 변형과 공유 작업 프레임}\label{sec:hra}
\textbf{오른팔 전용 계약} (\code{right_only.py}). 텐서 레이아웃은 바뀌지 않는다. 왼쪽 상태는 그리퍼 0인 항등 자세, 왼쪽 액션은 0
동작이며, 에피소드별 20차원 loss mask가 오른팔의 위치와 회전(9차원; 그리퍼는 정지해 있어 마스킹)만 지도한다.

\textbf{로봇 카메라 렌더링}(robotcam v2). 로봇 카메라의 각 픽셀에 대해, 측정된 로봇 내부 파라미터로부터의 광선을 어안 모델에
투영하여 재표본한다(어안 시야 밖의 광선은 검게 둔다). 작은 초점 거리·주점 지터와 같은 열화 사슬을 적용한다. 레이블은 핸드-아이
변환을 통해 UMI TCP에서 로봇 TCP로 옮긴다. $T_\text{robot}=T_\text{UMI}M^{-1}$.

\textbf{작업 프레임 G.} 테이프 선 위에서 로봇 조 끝으로 세 번 터치하여 원점(십자 중심), 테이프를 따르는 $x$, 베이스 $z$를 따르는
$z$를 얻는다. 인간 측에서 원점 자세는 정지 구간의 위치 중앙값과 평균 회전이다. 그 yaw는 테이크마다 십자 주변의 테이블 특징을 기준
테이크와 매칭하여(ORB, ratio test, PnP-RANSAC) 측정하고, 기울기는 IMU 중력(수정된 외부 파라미터)에서 얻는다. 상태는 G에서의 절대
자세이며 레이블은 $T_G = T_{\text{base},G}\,R(\text{yaw})\,T_\text{origin}^{-1}T_\text{robot}$이다.

\textbf{테이블 인지 변환.} 인간 궤적을 자신의 시작점에서 선택한 로봇 시작 자세 $T_S$ 쪽으로, 정규화 시간 $s$에 대한 smoothstep
가중치 $\mathrm{sst}(x)=3x^2-2x^3$으로 변형한다. $w_{xy}=1-\mathrm{sst}((s-0.2)/0.6)$, $w_z=1-\mathrm{sst}(s/0.3)$,
$w_\text{rot}=1-\mathrm{sst}((s-0.1)/0.5)$이므로 큐브 근처의 끝점은 인간의 것이 유지된다. 그다음 높이 보정 $c$를 테이크마다 이차
계획 문제로 푼다.
\begin{equation}
\begin{aligned}
\min_c\;& \textstyle\sum c^2+50(\Delta c)^2+2000(\Delta^2 c)^2+60\,m\,(c-g)^2\\
\text{s.t.}\;& z_\text{tip}+c\ge z_\text{table}+5\,\text{mm},
\end{aligned}
\end{equation}
여기서 $c_0=0$, 끝점 들어올림은 최대 30\mm{}이며, $g$는 낮은 프레임을 12\mm{} 간격 쪽으로 끌어올린다. 결과는 IK로 재생한다(두 번째
프레임마다; 오차 5\mm{}와 3\textdegree{} 미만, 관절 여유 2\textdegree{} 초과). 테이크는 PASS(들어올림 $\le$20\mm), PASS-CORRECTED
(20--30\mm), REVIEW, EXCLUDED, REJECT(IK 성공률 98\% 미만, 원점 정지 또는 yaw 누락, 정지 중 10\mm{} 넘게 이동, 바닥 아래에서
시작)로 분류하며, 앞의 두 부류만 내보낸다.

\section{모델}\label{sec:model}
LeRobot~\cite{cadene2024lerobot}을 통한 X-VLA~\cite{zheng2025xvla}. Florence-2~\cite{xiao2024florence2} 비전-언어 인코더, 24층
soft prompt transformer(hidden 1024, 헤드 16개), flow matching~\cite{lipman2023flow} 액션 헤드로 구성되며, 파라미터는 약 0.9B이고
모두 학습 가능하다. 입력은 패딩과 함께 224$\times$224로 크기를 맞춘 뷰 세 개, 지시문(BART 토크나이저), 고유수용 상태이다.
\textbf{Soft prompt.} 30개 도메인 $\times$ 32 토큰의 임베딩 테이블을 시퀀스에 덧붙이며, 상태/액션 프로젝션은 도메인별 선형층이다.
도메인 id는 체크포인트의 전처리기에 설정한다. 공개 가중치를 읽어 보니 도메인 10--17만 학습된 적이 있었으므로, 우리 embodiment에는
쓰이지 않은 도메인 20을 사용한다. \textbf{Flow matching.} 학습에서는 층화된 $t$를 뽑아 $x_t=\varepsilon t+a(1-t)$를 만들고 깨끗한
액션 $a$를 직접 회귀한다. 추론은 $x_1\sim\mathcal N(0,I)$에서 시작하여 $i=N\dots1$에 대해 $t=i/N$, $x_t=x_1 t+\hat a(1-t)$,
$\hat a=f(x_t)$로 두며 $x_1$을 재사용한다(기본 10스텝). \textbf{Gen-2를 위한 확장.} 공개 헤드(20차원 액션과 상태)를 32차원 액션과
20차원 상태로 넓혔다. 기존 행은 그대로 복사하고, 새 행은 기존 행의 스케일로 뽑으며, 새 바이어스는 0으로 두고, 확장된 모델이 모든
도메인에서 기존 모델을 재현하는지 확인하였다.

\section{학습}\label{sec:train}
\textbf{Gen-1} (\code{training/train_bi.py}). \ref{sec:gen1}절의 관절 변위 목표와 손실, AdamW(학습률 $10^{-4}$, 비전-언어 그룹은
0.1$\times$, weight decay $10^{-4}$, 그래디언트 클립 10), 1k 워밍업 후 $2.5\times10^{-6}$까지 코사인 감쇠, 배치 4, TF32를 쓰는 fp32,
300k 스텝, 10k마다 체크포인트. 먼저 20스텝 스모크 실행에서 유한한 손실과 활성화된 보조 항이 기록되어야 한다. 체크포인트 선택 규칙
(검증 geo 점수 최소)은 학습 전에 고정하였다.

\textbf{Gen-2, REL-only} (\code{training/relonly/}). 손실은 액션 차원 0--19에 대한 MSE이며, 학습과 추론 모두에서 모든 디노이징
스텝마다 차원 20--31이 입력과 출력에서 0이 되도록 transformer를 패치한다. 사전 점검(preflight)은 손실이 무작위화한 패딩 목표를
무시하는지, 패딩 예측이 정확히 0인지, 그 디코더 그래디언트가 사라지는지를 확인한다. 저장 shim을 둔 8-bit AdamW(bitsandbytes),
fp32 마스터 가중치를 둔 bf16 autocast, 배치 8, 시드 1000, 워밍업 1k, 정확히 마지막 스텝에서 끝나는 코사인 감쇠(LeRobot이 스케줄을
조용히 재조정하므로, 런처는 스텝 수와 감쇠가 다른 실행을 거부한다). 액션은 평균/표준편차로 정규화하고 상태는 정규화하지 않는다.
\textbf{Loss mask}(단일 팔): 마스킹된 목표는 뺄셈 전에 detach한 예측으로 대체하며(NaN 그래디언트 방지), 손실은 지도되는 원소 수로
정규화한다. \textbf{파인튜닝}은 사전학습 가중치에서 새 옵티마이저로 시작하며 로봇 데이터셋의 정규화 통계를 사용한다. 에피소드
부분집합(예: R30)은 인덱스 목록으로 선택한다. \textbf{공동 학습}은 배치마다 정확히 에고 샘플 네 개와 로봇 샘플 네 개를 뽑는다.

\textbf{운영.} 작업은 Ray 클러스터에 제출하여 GPU 노드(RTX~4090 또는 RTX~5090)에 고정한다. 두 노드의 PyTorch 버전이 달라 노드 간
실행은 비트 단위로 같지 않다. 런처는 사용 중인 GPU, 코드 해시나 도메인 id 불일치, 그리퍼 계약 게이트 실패, 디스크 공간 부족 시
시작을 거부하며, 디스크 가드는 노드가 가득 차기 전에 학습을 멈춘다. 각 본 실행 전에 60스텝 스모크 실행(저장·재개 포함)을 한다.
아카이버는 모든 체크포인트를 SHA-256 검증과 함께 외부 저장소로 복사한 뒤 노드에서 삭제하며, 옵티마이저 상태는 고정 간격으로만
보존한다.

\section{평가}\label{sec:eval}
\textbf{Gen-1 오프라인} (\code{evaluation/c8old_mac_eval.py}). held-out 에피소드에서 세그먼트당 시작 프레임 일곱 개(0--30에서 5프레임
간격)를 쓰고, 시작마다 고정 시드로 flow 샘플 하나를 뽑는다. 호라이즌 $k\in\{1,4,8,16,30\}$에 대해 TCP 오차
$\lVert\mathrm{FK}(q_t+\hat a_k)-p_k\rVert$(중앙값, p90), 관절 MAE, 0 동작 예측기의 MAE, 예측과 측정 TCP 변위 사이의 코사인을
계산한다. geo 점수는 호라이즌별 TCP 오차 중앙값의 평균이며, collapse 비율은 예측과 목표의 퍼짐을 비교한다. 움직이는 부분집합은
$k{=}30$까지 TCP가 20\mm{} 이상 움직인 샘플이다. \textbf{지시문 교체}: 각 샘플을 한 시드에서 여섯 순서 지시문 모두로, 그리고 참
지시문으로 시드 두 개를 더 써서 실행한다. 지시문에 따른 퍼짐을 시드에 따른 퍼짐과 비교하고, 여섯 개 중 참 지시문의 오차 순위를
우연 수준과 비교한다. \textbf{대응 부트스트랩}: 두 초기화 사이의 geo 점수 상대 차이를 에피소드 수준 재표본(10k회, 두 모델에 동일)으로
추정한다.

\textbf{Gen-2 진단} (\code{analysis/v3/}). 로봇 프레임에서의 teacher-forced 방향 코사인($k$=4/8/16에서 10\mm{} 넘게 움직이는 프레임에
대한 예측 대 기록 병진)을 모든 좌우 교체와 축 부호 반전에 대해 반복한다. 이미지/상태 도메인 교체(에고 또는 로봇 이미지를 에고,
로봇 또는 0 상태와 결합), 손목 비디오의 optical flow 대 자세 축 상관(컨벤션 검사), 라플라시안 분산으로 잰 손목 이미지 선명도도
본다. \textbf{학습 손실 비교}: 손실을 로그에서 파싱하고(스텝 = 줄 인덱스 $\times$ 로그 간격) 창 평균과 비로 비교한다.
\textbf{Held-out 손실}(단일 팔): 모든 검증 행에 대해 노이즈 시드 두 개로, 그리고 2{,}000행 학습 부분집합에 대해 (마스킹된) 학습
손실을 계산한다.

\textbf{주장 검증} (\code{analysis/verify_claims.py}). 보고서의 모든 수치를 저장소의 증거 파일로부터 인쇄된 정밀도로 다시 계산하며,
보고서 소스와 README에서 필수 문자열과 폐기된 문자열을 검사한다.

\section{로봇 배포}\label{sec:deploy}
\textbf{추론.} 정책은 Mac(Apple 실리콘 GPU, bf16)에서 디노이징 스텝 5회(양팔) 또는 10회(단일 팔)와 REL-only 0 패치로 실행된다. 로봇
카메라 세 대를 640$\times$480으로 캡처하고, 선택적으로 확대한 뒤(오른 손목, 단일 팔) 224$\times$224로 크기를 맞추어 학습과 같이
정규화한다. 단일 팔 모드에서는 왼 손목 이미지가 검은색이다. \textbf{상태}는 학습 계약을 따른다. 측정 관절의 FK로 구한 TCP 자세에
FK와 데이터셋 TCP 컨벤션 사이의 고정 프레임 보정을 적용하고, 실행 시작 시 잡은 앵커 또는 고정 앵커 파일(원점 고정 모델은 G) 기준
상대 자세로 표현한다. \textbf{디코딩}: $k=1..16$에 대해 절대 목표 $T_k=T_\text{now}A_k$. 그리퍼 값은 0--45 명령으로 대응되며(로봇
raw 0 = 닫힘, $-270$ = 열림), 단일 팔 모드에서는 조를 고정한다.

\textbf{IK}(Pink~\cite{caron2024pink}). 손가락을 잠근 URDF를 쓰며, TCP는 그리퍼 링크에서 $x$를 따라 70\mm{} 떨어져 있다(데이터셋
TCP와 동일). 팔별 프레임 태스크(위치 1.0, 방향 0.5)와 가벼운 자세 태스크를 측정 관절에서 이차 계획법으로 풀며, 관절 한계 여유와
웨이포인트별 관절 스텝 한계를 둔다. 프로파일에 따라 손목 yaw 관절을 잠그거나 풀며, 다른 솔버(수치, 학습 기반, GPU 배치)로 실행 중에
전환할 수 있다.

\textbf{실행.} \emph{스텝 모드}는 각 청크의 웨이포인트 몇 개를 실행하고 팔이 안정될 때까지 기다린다. \emph{PLAN mode}는 순전파 하나를
항상 진행 중으로 유지하면서 IK로 푼 행(50.05\,ms 간격)을 로봇 서비스로 스트리밍하고, 로봇 서비스는 이를 속도 피드포워드와 함께
100\,Hz 보간점(knot)으로 재생한다. 예상 지연 시간 안에 실행될 행은 예약해 두고 그 뒤의 행만 교체한다. 청크 사이 이음매는 관절 공간 지수
블렌드(\code{plan_ema}), 선택적 크로스페이드, 선택적 다항식 평활화, 다음 순전파 전 최소 실행 행 수로 처리한다. 이 값들은 동작을 눈에
띄게 바꾸므로 실행마다 기록한다. \emph{Real-time chunking}(선택): 각 flow 스텝에서 예측을 이전 청크의 아직 실행되지 않은 목표 쪽으로
끌어당긴다. $\hat x\leftarrow\hat x+t\,g_w\,J^\top(W\odot(Y-\hat x))$이며, 가이던스 가중치는 10으로 상한을 두고 $W$는 처음 8스텝에
걸쳐 감쇠한다~\cite{black2025rtc}.

\textbf{로봇 서비스.} 적분 피드포워드 항이 있는 MIT impedance mode의 팔 컨트롤러 두 개, 100\,Hz 제어 스레드, 가속도 제한 동기화
프로파일로 구성된다. raw 관절 한계는 클리핑하며, 계획된 보간점이 90\textdegree{} 넘게 튀면 팔을 정지시킨다. \textbf{안전과 정지.} NaN
예측이면 중단한다. 단일 팔 모드에서는 왼팔과 조가 항상 고정된다. 이동량, 스텝, 출력 가드는 선택 사항이다. \emph{큐브 크기 정지}는
raw 손목 프레임의 HSV 마스크로 빨간 큐브의 겉보기 크기(픽셀 단위 $\sqrt{\text{area}}$)를 0.1\,s마다 재고, 임계값을 넘는 판독이 두 번
나오면 멈춘다. 선택적인 추종 모드는 크기가 임계값의 85\% 아래로 줄면 재개한다. \textbf{체크포인트}는 학습 노드나 아카이브에서
가져와 SHA-256과 신원(실행, 데이터셋, 도메인)으로 검증하고, 인터페이스를 재시작하여 로드한다. \textbf{시행 기록}: 작업자 양식에
체크포인트, 모델 종류, 로봇 데이터, 지시문, 배치, 쌓기 단계 0--3, 잘못된 색, 실패 사유, 실행의 정지 사유를 기록하며, 성공은 색이
맞는 단계 3으로 정의한다.

\section{알려진 불일치}\label{sec:issues}
이 문서를 작성하면서 다음 사항이 드러났으며, 숨기지 않고 나열한다. (1) 팔 간 간격 검사는 존재하지만 Gen-1 게이트가 아니라 리타게팅
v2에서만 쓰인다(이전 README 문장이 다르게 적고 있었으며 수정하였다). (2) Gen-1 목표 코드에서 이진화된 그리퍼 레이블 1은 주석에는
닫힘이라고 되어 있지만 실제로는 \emph{열림}을 뜻한다. (3) 그리퍼 캘리브레이션이 세 가지 사용된다(67.5\mm{} 개구의 Gen-1 tick, Gen-2의
gcal1 캘리퍼 보간점, 수집기의 열림/닫힘 파일). 각 데이터셋은 자신이 쓰는 것을 명시한다. (4) 내보내기 도구가 쓰는 카메라--IMU 시간
오프셋은 raw 캘리브레이션 출력과 부호와 크기가 다르며, 이를 변환한 단계는 기록되어 있지 않다. (5) 내보내기 도구와 한 뷰어 도구가 서로
다른 두 버전의 오른쪽 어안 캘리브레이션을 사용한다. (6) IK 스텝 임계값은 15\,Hz 행 기준으로 문서화되어 있으나 30\,fps 행에 적용된다.
(7) PLAN mode에서 양팔 경로는 그리퍼 폭을 웨이포인트 솔버에 전달하지 않으므로 조가 닫힌 채로 남을 가능성이 높다. 조를 고정하는
단일 팔 실행은 영향을 받지 않는다. (8) 여러 진단은 저장소 밖의 스크립트로 실행되었으며, 그 출력은 다시 실행한 것이 아니라 캡처한
로그로 보관되어 있다. (9) G 프레임, 손가락 끝 박스, 18\textdegree{} 카메라--IMU 보정의 피팅 스크립트는 저장소에 없고 그 출력만 있다.

\begin{thebibliography}{99}\footnotesize
\bibitem{chi2024umi} C.~Chi \emph{et al.}, ``Universal manipulation interface: In-the-wild robot teaching without in-the-wild robots,'' in \emph{RSS}, 2024.
\bibitem{furgale2013kalibr} P.~Furgale, J.~Rehder, and R.~Siegwart, ``Unified temporal and spatial calibration for multi-sensor systems,'' in \emph{IROS}, 2013.
\bibitem{tsai1989handeye} R.~Y. Tsai and R.~K. Lenz, ``A new technique for fully autonomous and efficient 3D robotics hand/eye calibration,'' \emph{IEEE Trans. Robotics and Automation}, 1989.
\bibitem{murai2025mast3rslam} R.~Murai, E.~Dexheimer, and A.~J. Davison, ``MASt3R-SLAM: Real-time dense SLAM with 3D reconstruction priors,'' in \emph{CVPR}, 2025.
\bibitem{zheng2025xvla} J.~Zheng \emph{et al.}, ``X-VLA: Soft-prompted transformer as scalable cross-embodiment vision-language-action model,'' in \emph{ICLR}, 2026.
\bibitem{cadene2024lerobot} R.~Cadene \emph{et al.}, ``LeRobot: State-of-the-art machine learning for real-world robotics in PyTorch,'' 2024.
\bibitem{xiao2024florence2} B.~Xiao \emph{et al.}, ``Florence-2: Advancing a unified representation for a variety of vision tasks,'' in \emph{CVPR}, 2024.
\bibitem{lipman2023flow} Y.~Lipman \emph{et al.}, ``Flow matching for generative modeling,'' in \emph{ICLR}, 2023.
\bibitem{caron2024pink} S.~Caron \emph{et al.}, ``Pink: Python inverse kinematics based on Pinocchio,'' \url{https://github.com/stephane-caron/pink}, 2024.
\bibitem{black2025rtc} K.~Black, M.~Y. Galliker, and S.~Levine, ``Real-time execution of action chunking flow policies,'' \emph{arXiv:2506.07339}, 2025.
\end{thebibliography}
\end{document}
